\section{System Model}
\label{sec:system}

\paragraph{Collaborative perception setting.}
An ego agent and a cooperating agent (CAV) each run their own object detector over their own sensor data and hold a set of locally detected boxes in their own coordinate frame. The CAV transmits its locally detected boxes together with its self-reported pose, its belief about its own position and orientation in the frame the two agents share. The ego agent uses that pose to project the CAV's boxes into its own frame before fusing them with its own detections. We describe one ego-CAV pair. With more than one cooperating agent, the ego agent repeats this for each.

\paragraph{Agent and coordinate frames.}
Each agent's locally detected boxes are expressed in that agent's own sensor frame. The ego agent's detections stay in the ego frame throughout and are written $\mathcal{P}=\{(\mathbf{p}_i,\psi_i,s_i)\}_{i=1}^{M}$. The CAV's locally detected boxes are transformed into the ego frame using the CAV's self-reported pose, giving $\mathcal{Q}=\{(\mathbf{q}_j,\phi_j,s_j)\}_{j=1}^{N}$. $\mathbf{p}_i,\mathbf{q}_j\in\R^2$ are BEV box centres and $\psi_i,\phi_j$ the corresponding headings, both already expressed in the ego frame. $s_i,s_j$ are the detector's confidence scores.

\paragraph{Shared bounding-box representation.}
Each transmitted detection is a box centre, extents and heading, seven floats, together with its detection confidence score. No appearance feature, learned embedding or map is part of the message. A received detection is exactly that seven-float-plus-score tuple, plus the sender's self-reported pose used to place it in the ego frame.

\paragraph{Misalignment model.}
The CAV's self-reported pose differs from its true pose. GNSS and odometry report position and orientation with an error that neither agent knows at inference time. Projecting the CAV's boxes with the reported pose instead of the true one applies the same rigid error to the whole box set at once, since every box of that agent is transformed by the same pose. We model this as an unknown rigid transform $T=(R(\theta),\mathbf{t})\in\SE$ such that $T\mathbf{q}_j$, not $\mathbf{q}_j$, is where the $j$-th object actually is. $T$ is planar. It captures a translation and a yaw error and leaves roll, pitch and height to the detector. Fig.~\ref{fig:misalignment} illustrates $T$ in the shared frame.

\begin{figure}[t]
  \centering
  \resizebox{\linewidth}{!}{% Fig: the misalignment model. Geometric, algorithm-independent (no AlignFormer).
\begin{tikzpicture}[
  >=Latex, scale=0.78, every node/.style={transform shape},
  obj/.style={fill=black!25, draw=black!50, rounded corners=0.5pt},
  egomk/.style={draw=green!45!black, line width=0.9pt, fill=green!45!black, fill opacity=0.15, rounded corners=0.5pt},
  cavmk/.style={draw=orange!70!black, line width=0.9pt, fill=orange!70!black, fill opacity=0.15, rounded corners=0.5pt},
  tick/.style={line width=0.8pt},
]
  % road band
  \fill[black!5] (-0.6,0.2) rectangle (9.2,4.4);
  \draw[black!25, dashed] (-0.6,2.3) -- (9.2,2.3);

  % three true object positions (gray silhouettes), rotated slightly for variety
  \foreach \x/\y/\a in {1.6/3.1/8, 4.4/1.3/-6, 7.2/3.4/4} {
    \begin{scope}[shift={(\x,\y)}, rotate=\a]
      \node[obj, minimum width=0.62cm, minimum height=0.34cm] {};
    \end{scope}
  }

  % ego boxes p_i: tight, correct, on every object
  \foreach \x/\y/\a in {1.6/3.1/8, 4.4/1.3/-6, 7.2/3.4/4} {
    \begin{scope}[shift={(\x,\y)}, rotate=\a]
      \node[egomk, minimum width=0.78cm, minimum height=0.46cm] {};
      \draw[tick, green!45!black] (0,0) -- (0,0.33);
    \end{scope}
  }

  % CAV boxes q_j: same uniform offset + rotation on every object (the rigid error)
  \foreach \x/\y/\a in {1.6/3.1/8, 4.4/1.3/-6, 7.2/3.4/4} {
    \begin{scope}[shift={($(\x,\y)+(0.62,0.5)$)}, rotate={\a+16}]
      \node[cavmk, minimum width=0.78cm, minimum height=0.46cm] {};
      \draw[tick, orange!70!black] (0,0) -- (0,0.33);
    \end{scope}
  }

  % highlight T on the middle object
  \coordinate (qmid) at ($(4.4,1.3)+(0.62,0.5)$);
  \coordinate (pmid) at (4.4,1.3);
  \draw[->, line width=0.8pt, black!70, dashed] (qmid) to[bend left=18] node[midway, above, font=\scriptsize, fill=white, inner sep=1pt] {$T=(R(\theta),\mathbf{t})$} (pmid);

  % legend
  \node[egomk, minimum width=0.3cm, minimum height=0.22cm] at (0.0,4.75) {};
  \node[font=\scriptsize, anchor=west] at (0.22,4.75) {$\mathcal{P}$ (ego, true position)};
  \node[cavmk, minimum width=0.3cm, minimum height=0.22cm] at (4.6,4.75) {};
  \node[font=\scriptsize, anchor=west] at (4.82,4.75) {$\mathcal{Q}$ (CAV, reported pose)};

  % axis icon: one shared ego frame
  \draw[->, line width=0.7pt] (-0.4,0.6) -- (0.35,0.6) node[right, font=\scriptsize] {$x$};
  \draw[->, line width=0.7pt] (-0.4,0.6) -- (-0.4,1.35) node[above, font=\scriptsize] {$y$};
  \node[font=\scriptsize, anchor=west] at (-0.4,-0.05) {ego frame (shared)};
\end{tikzpicture}
}
  \caption{\textbf{The misalignment transform $T$.} Grey silhouettes are the true object positions. $\mathcal{P}$ (green) is the ego's own detections, placed correctly. $\mathcal{Q}$ (orange) is the CAV's detections as received under its self-reported pose, displaced from every object by the same rigid transform $T$. Both are expressed in the one shared ego frame.}
  \label{fig:misalignment}
\end{figure}

\paragraph{Problem formulation.}
The ego agent observes $\mathcal{P}$ and $\mathcal{Q}$ only. It has no access to the CAV's true pose and no given correspondence between its own and the CAV's detections of the same physical object. The problem is to recover a correction close to $T$ from $\mathcal{P}$ and $\mathcal{Q}$ alone and apply it to $\mathcal{Q}$, so the corrected boxes can be handed to an otherwise unchanged fusion stage together with $\mathcal{P}$. The system is scored by the detection accuracy of the fused output, not by how close the recovered correction is to $T$. Sec.~\ref{sec:protocol} shows why that distinction is not academic.
